\begin{statement}{E.41}{Proved}{The complete two-adic class-two fibre and its exact finite lifting images}
Put
\[
 P_{2^n}=B/\bigl(\gamma_3(H)H^{2^n}\bigr),\qquad
 \mathscr J_{2^n}=\operatorname{im}
       \bigl(B\longrightarrow\mathscr D\times P_{2^n}\bigr).
 \tag{E.41.1}
\]
Fix any one of the $432$ multiplier-one source words
of E.24.  At exponent four its complete correction-image
fibre has $2^{10}$ elements, and every one lifts through
the whole two-adic class-two tower.  For $n\ge3$ the
complete fibre at $P_{2^n}$ has
\[
                 2^{10n-9}
 \tag{E.41.2}
\]
elements.  Exactly $2^{10n-10}$ lift to the next level,
and this is also exactly the subset lifting through
every level.  The criterion is the single additional
congruence in \textup{(E.41.22)} below.

The full inverse-limit fibre has an explicit
homeomorphism with $\mathbb Z_2^{10}$.
Its affine source congruences, integral Hall parity
and eight central parameters are retained in
\textup{(E.41.15)--(E.41.20)}.
It can be prescribed simultaneously with any full
three-adic fibre from E.37, by one profinite
commutator word.  Every finite common shadow so
obtained is generating.

There is an individually specified exponent-eight
shadow over source row $0$, parameters $(0,0;0)$,
which has no multiplier-one commutator lift to
$P_{16}$.  Its actual source-word construction and
Hall coordinates are given in
\textup{(E.41.23)--(E.41.26)}.
\end{statement}
\begin{proof}
We first make the integral coordinate comparison.
Let the ten ordered Schreier generators of $H$
be the ones fixed in E.33.
For Hall coordinates $(v,z)$, with central basis
$[\eta_i,\eta_j]$ ordered by $i<j$, define
\[
       Z_{ij}=2z_{ij}+v_i v_j.
 \tag{E.41.3}
\]
This is an injective map into doubled-log
coordinates, with image exactly
$Z_{ij}\equiv v_i v_j\pmod2$.
Its inverse on that image is
$z_{ij}=(Z_{ij}-v_i v_j)/2$.
The actual multiplication becomes
\[
             (v,Z)(w,T)=(v+w,Z+T+v\wedge w).
 \tag{E.41.4}
\]
The signed Schreier substitutions for $s,r$
give the same integral matrices $S,R$ as E.37,
with $S^2=R^3=1$ over $\mathbb Z$.
In particular they preserve the displayed
parity subset because they are collected from
actual free-group words.

In the exponent-$2^n$ class-two quotient, Hall
abelian coordinates are modulo $2^n$ and Hall
central coordinates are modulo $2^{n-1}$.
Indeed $[x,y]^{2^n}=1$ follows from commutator
collection, and the power relation for $(xy)^{2^n}$
also gives $[x,y]^{\binom{2^n}{2}}=1$;
their two-adic greatest common divisor is $2^{n-1}$.
Conversely the Hall group with exactly these
coordinate moduli satisfies every $2^n$-th power
relation, by the same power formula.
These two actual homomorphisms are inverse.
Equivalently, both $v,Z$ may be reduced modulo
$2^n$ while retaining \textup{(E.41.3)}.
Thus the Hall central precision is one power of
two below the abelian precision.

\paragraph{The complete source image and actual word preimages.}
Write $R_2$ for the kernel on
$K_6\cap\widehat F'$ of the full binary source
quotient in E.24.  Its $5458$ normal relators
were retained in E.36, including the $95$
primitive integral-$F$-abelianization relations.
Their integral $H$-abelian image is precisely
the eight-column lattice $K$ of E.31.
The new integral collection checks all $5458$
words against their previous exponent-four
values.  Twelve actual relators, with indices
\[
       3,4,96,97,98,100,103,104,105,4656,4658,4662,
 \tag{E.41.5}
\]
give ordered-product words $g_1,\ldots,g_8$
whose abelian vectors are exactly the columns
of $K$.  The saved integer combination matrix
has largest absolute coefficient $68$.

For each source relator $\rho$, replace it by
$\rho(\prod_i g_i^{c_i})^{-1}$, where
$[\rho]_H=Kc$.  Also retain the similarly adjusted
$s,r$ conjugates of the eight $g_i$ and all
$28$ commutators $[g_i,g_j]$.
All $5502$ resulting words have zero integral
$H$-abelianization.
Their central images, closed under the actual
exterior $s,r$ action, form a lattice
$C\subset\mathbb Z^{45}$ with
\[
             \operatorname{Smith}(C)=(1^{39},2^4,4^2).
 \tag{E.41.6}
\]
This is the full central image.
To check completeness, quotient the normal
source presentation by these adjusted words.
Every source relator becomes an ordered product
of the $g_i$; their conjugates remain such products,
and the commutator relations make that quotient
abelian.  Their eight independent integral
abelian vectors then identify the quotient with
$K$.  Hence the kernel of its abelian map is
exactly the normal closure just collected.
In the class-two image, conjugation acts on
this kernel by the displayed exterior matrices,
which is precisely the closure used to compute
$C$.

This calculation also retains actual central
word preimages.  Forty-five selected adjusted
relators $w_1,\ldots,w_{45}$ already have central
matrix $W$ with
\[
 W=C A,\qquad
 A\in M_{45}(\mathbb Z),\qquad
 |\det A|=631022156534756593832815423865613.
 \tag{E.41.7}
\]
The determinant is odd.  Therefore these literal
zero-$H$-abelianization words give a basis of
$C\otimes\mathbb Z_2$.
The certificate specifies every adjusted
relator and its ordered $g_i$ correction,
and checks its integral central vector.
It does not infer word preimages merely from
the rational rank.

The same two-adic image is obtained from the
kernel of the entire source $\mathscr D$.
Every word used above lies in $K_6$, whose
ternary source image has exponent three.
Cubing clears that image, and multiplication
by three is invertible in $\mathbb Z_2$.
Equivalently use continuous two-supported
exponents on these actual words.
Thus the complete two-adic source image is
described as follows: choose $c\in\mathbb Z_2^8$,
form the actual ordered product $\prod_i g_i^{c_i}$
in the class-two pro-two group, and then multiply
by an arbitrary central element of
$C\otimes\mathbb Z_2$ in Hall coordinates.
This is a closed image, since it is the image
of the compact coefficient space by the
displayed continuous word map.

We also need the exact finite fibre over an
abelian value.  It is the reduction of $C$,
not a larger uncomputed kernel.
In the chosen $K$ basis, $g_1,g_2$ have
abelian vectors $12e_1,12e_2$.
The remaining six columns have a unit triangular
minor.  Therefore the coefficient ambiguities
modulo $2^n$ are generated by $2^{n-2}g_1$,
$2^{n-2}g_2$ and $2^n$ multiples of the other
coordinates.  The Hall central values of
$g_1,g_2$ belong to $C+2\mathbb Z^{45}$,
as the saved literal parity calculation verifies.
Their $2^{n-2}$ powers consequently have
central values in $C$ modulo $2^{n-1}$.
The $2^n$ powers and all collecting terms involving
an abelian multiple of $2^n$ vanish at that
central precision, by the Hall power formula.
This proves the finite comparison for every
$n\ge2$.

\paragraph{The full central norm and its finite kernels.}
Let $M$ be the $90$-by-$45$ central norm matrix
from E.37; its saved integral decomposition is
\[
 S_M=U_M M V_M,\qquad
 \operatorname{diag}(S_M)=(1^{35},3,3,0^8).
 \tag{E.41.8}
\]
It is saturated over $\mathbb Z_2$.
The complete source norm has
\[
             \operatorname{Smith}(MC)
                      =(1^{31},2^4,12^2,0^8).
 \tag{E.41.9}
\]
Put $N_0=(V_M)_{\{37,\ldots,44\}}$,
the last eight columns, with zero-based indexing.
The exact certificate verifies
\[
           M N_0=0,\qquad C^{-1}N_0\in M_{45,8}(\mathbb Z).
 \tag{E.41.10}
\]
The columns of $N_0$ are primitive in the
ambient central lattice, since $V_M$ is
unimodular.  They therefore give all integral
central norm-kernel parameters inside the
source lattice.

For central modulus $2^m$, the Smith factors
give respectively
\[
 \begin{split}
 |C\bmod2^m|&=
  2^{39m+4\max(m-1,0)+2\max(m-2,0)},\\
 |MC\bmod2^m|&=
  2^{31m+4\max(m-1,0)+2\max(m-2,0)}.
 \end{split}
 \tag{E.41.11}
\]
Thus the kernel on correction images has
$2^{8m}$ elements.
The contained image $N_0(\mathbb Z/2^m)^8$
has that same order by primitivity.
They are equal.  In particular every finite
homogeneous central correction image lifts
compatibly through all levels.

Here is an exact small congruence map for the
remaining affine source obstruction.
If $\sigma_i=(S_M)_{ii}$ for $i<37$, define
\[
 \begin{split}
 A_M&=-3(V_M)_{<37}
       \operatorname{diag}(\sigma_i^{-1})
       (U_M)_{<37},\\
 D_C&=4C^{-1}A_M .
 \end{split}
 \tag{E.41.12}
\]
Both matrices are integral, as verified entry
by entry.  On the saturated ambient norm image,
$A_Mh/3$ solves $Mz+h=0$.
Since three is a two-adic unit, that solution
lies in $C\otimes\mathbb Z_2$ exactly when
$D_Ch=0$ modulo four.
All ambient homogeneous solutions differ by
$N_0$, already contained in $C$, so this is
the complete source condition.
The identity $D_C MC=0\pmod4$ is checked
integrally.  The same argument works modulo
$2^m$ for $m\ge2$: an ambient solution and this
modulo-four condition give a source solution
modulo $2^m$.  Exponent four was computed
directly in E.36.

\paragraph{The source conic and all $432$ origins.}
The exact abelian norm kernel is the primitive
basis from E.37:
\[
 Bt=(-t_1,t_2,-t_1+t_3,t_2,t_3,-t_3,
             t_1-t_3,t_3,-t_1,t_1)^t.
 \tag{E.41.13}
\]
Its defining equations, including $Lv=0$,
have seven unit Smith factors over $\mathbb Z$.
They therefore have this same kernel modulo
every $2^n$.
For the two source functionals $\lambda_1,\lambda_2$
in \textup{(E.37.6)}, the full two-adic
$K$-membership condition is divisibility by
four of both.  Their common kernel has the
same index $2^4$ as $K$ in $\ker L$.
Set
\[
 u=t_1+t_2,\qquad v=t_3-t_2,\qquad w=t_2,
 \qquad
 (t_1,t_2,t_3)=(u-w,w,v+w).
 \tag{E.41.14}
\]
On $Bt$ the two functionals are
$10u-v-12w$ and $u-v-12w$.
For the retained source word $f$, write
$c_i=\lambda_i([f])\pmod4$.
The exact affine condition is therefore
\[
        u\equiv c_1-c_2=:u_0\pmod4,\qquad
        v\equiv c_1-2c_2=:v_0\pmod4.
 \tag{E.41.15}
\]
Every integral source word is collected and
compared with both the previous modulo-three
BCH value and modulo-four Hall value.
The resulting six residue pairs are
\[
        (0,0),\ (2,0),\ (0,1),\
        (2,3),\ (3,1),\ (3,3),
 \tag{E.41.16}
\]
each occurring $72$ times.
For CM rows $0,\ldots,23$ the pairs, respectively,
are the entries of the following array, read
across rows; each is constant on all eighteen
internal choices:
\[
 \begin{array}{cccc}
 (0,0)&(2,0)&(2,0)&(0,0)\\
 (0,1)&(2,3)&(2,3)&(0,1)\\
 (3,1)&(3,3)&(3,3)&(3,1)\\
 (3,1)&(3,3)&(3,3)&(3,1)\\
 (0,0)&(2,0)&(2,0)&(0,0)\\
 (0,1)&(2,3)&(2,3)&(0,1)
 \end{array}
 \tag{E.41.17}
\]

The rational-cokernel polynomial of E.37 becomes
\[
                 F(u,v)=v+3(u^2+uv+v^2).
 \tag{E.41.18}
\]
Its doubled-log zero-Smith rows are even
integer multiples of $F$, with one exactly
$2F$ and greatest common divisor of their
multipliers equal to two.
For an actual source reference $g_0$ with
abelian value $Bt$, both norm words have
zero integral abelianization and even doubled
central values.  Let $h$ be those central
values divided by two.
The Hall zero-Smith rows are consequently
integer multiples of $F$ with greatest common
divisor one.  The exact ambient obstruction
is $F=0$; at exponent $2^n$ it is
$F=0$ modulo $2^{n-1}$.

The remaining source congruence
$D_Ch=0\pmod4$ is automatic on all the affine
source lattices \textup{(E.41.15)}.
This assertion is checked on the complete
finite source, as follows.
The homogeneous abelian source norm basis,
in the $t$ coordinates, is
\[
 \begin{pmatrix}
 1&0&0\\
 43&-12&4\\
 -33&0&-4
 \end{pmatrix}.
 \tag{E.41.19}
\]
Its first column is primitive and its last
two columns divided by four complete a
two-adic basis.  Hence, above each of the
$432$ actual origins, its coefficients
modulo $(8,2,2)$ enumerate exactly the
$32$ possible abelian norm values modulo eight.
Using the actual $g_i$ word lifts, the program
forms all $432\cdot32$ reference values,
computes their Hall errors $h$, and verifies
$D_Ch=0$ modulo four in every case.
It then constructs the source central
correction and checks the two norm equations
modulo eight.  All $32$ values pass over
every source point; all resulting geometric
Frattini ranks are ten.

This finite calculation proves the asserted
full congruence because its only modulus is
four.  The finite source comparison above
shows that any two actual reference lifts
with the same abelian value modulo eight
differ by a central element of $C$ modulo
four.  Their norm errors differ by $MC$,
annihilated by $D_C$ modulo four.
Thus the finite test covers every
two-adic reference value, including its
central carries.

We can now give the full coordinate map.
When $u_0$ is even, choose
$u=u_0+4a$ and let $v$ be the unique root
of $F(u,v)=0$ with $v\equiv v_0\pmod4$.
When $u_0$ is odd, choose $v=v_0+4a$
and solve uniquely for $u\equiv u_0\pmod4$.
The first case has
$F_v=1+3(u+2v)$ odd; the second has
$F_u=3(2u+v)$ odd.
All six initial pairs satisfy $F=0$ modulo
four, so the roots exist uniquely by the
explicit successive-bit Hensel construction.
Choose $w\in\mathbb Z_2$ arbitrarily and
recover $t$ by \textup{(E.41.14)}.
Solve $Kc=Bt-[f]$ and form
$g_0=f\prod_i g_i^{c_i}$.
Then all the central solutions, in Hall
correction coordinates, are
\[
       z_{\rm corr}=A_Mh/3+N_0b,\qquad
       b\in\mathbb Z_2^8.
 \tag{E.41.20}
\]
The congruence just proved puts this vector
in $C\otimes\mathbb Z_2$.
Actual central preimages are supplied by
$W$ in \textup{(E.41.7)}, using
$W^{-1}z_{\rm corr}\in\mathbb Z_2^{45}$.
All Hall parities are retained because the
construction starts from actual ordered
source words and adds an actual central
Hall vector.
The parameters $(a,w,b)$ range over
$\mathbb Z_2^{10}$.
Conversely recover $a,w$ from the target
abelian vector, subtract the displayed
particular correction, and recover $b$
using the inverse unimodular Smith matrix.
These continuous inverse maps prove the
claimed homeomorphism.

\paragraph{Exact finite images and the first failure.}
For $n\ge3$, conic solutions modulo $2^{n-1}$
in a fixed residue disk modulo four have
$2^{n-3}$ choices.  Both coordinates then
have a free top bit when reduced modulo
$2^n$, giving $2^{n-1}$ pairs $(u,v)$.
The independent $w$ gives $2^{2n-1}$
abelian targets.
Each has exactly $2^{8(n-1)}$ central
correction images by \textup{(E.41.11)}.
This proves \textup{(E.41.2)}.
At $n=2$ there are four abelian targets,
and the central count is $2^8$, as in E.36.

An abelian target modulo $2^n$ lifts to an
exact conic point precisely when
$F=0$ modulo $2^n$.
The same condition is necessary at the
next quotient, whose Hall central precision
is $2^n$, and sufficient by the unit
derivative in the same residue disk.
For $n\ge2$, the number of such abelian
targets is $2^{2n-2}$.
There is no additional central loss:
choose an exact conic lift and one exact
central solution; the discrepancy from
any prescribed finite central solution
belongs to the finite homogeneous kernel,
which is exactly the reduction of $N_0$
by \textup{(E.41.11)}.
Consequently
\[
 \#\{\text{images lifting through all levels}\}
                      =2^{10n-10}\qquad(n\ge2),
 \tag{E.41.21}
\]
and the full criterion is
\[
             F(u,v)\equiv0\pmod{2^n}.
 \tag{E.41.22}
\]
At exponent four it holds for every image;
at every subsequent level it selects
exactly one half.

Here is an explicit failing image at exponent
eight.  Use source row $0$, parameters $(0,0;0)$,
and the abelian target
\[
             v_H=(0,0,4,0,4,4,4,4,0,0)\pmod8.
 \tag{E.41.23}
\]
Its Hall central vector, in the ordered
$45$-coordinate wedge basis, is the
concatenation of these three rows:
\[
 \begin{array}{rrrrrrrrrrrrrrr}
 2&2&3&0&1&3&2&2&3&0&0&0&2&2&2\\
 1&0&3&1&1&3&3&2&1&1&3&2&0&1&2\\
 0&2&1&2&3&3&2&1&0&3&3&2&0&3&1
 \end{array}
 \pmod4 .
 \tag{E.41.24}
\]
It is an actual common-source correction.
If $f_0$ is the retained literal source word,
take
\[
          f_0\prod_{i=1}^8g_i^{\,9c_i}
                 \prod_{j=1}^{45}w_j^{\,9a_j},
 \qquad c=(0,3,5,1,5,3,7,1),
 \tag{E.41.25}
\]
where the $a_j$ are the concatenated rows
\[
 \begin{array}{rrrrrrrrrrrrrrr}
 2&0&0&2&3&1&3&3&1&3&1&2&1&3&3\\
 3&0&0&3&1&2&1&2&2&0&3&0&0&1&2\\
 3&0&3&3&1&2&2&3&0&0&0&1&3&3&2 .
 \end{array}
 \tag{E.41.26}
\]
The ninth powers clear every ternary source
coordinate and preserve the exponent-eight
image; all the words already belong to the
binary source kernel.
The literal coordinate check gives
\textup{(E.41.23)--(E.41.24)} and verifies
both norm relations.
Every lift to $P_{16}$ would retain
$(t_1,t_2,t_3)=(0,0,4)$ modulo eight,
so $(u,v)=(0,4)$ modulo eight.
But $F(0,4)=52\equiv4\pmod8$, contradicting
the necessary Hall zero-Smith row at $P_{16}$.
This is the asserted failure for the fixed
multiplier-one commutator shadow.

\paragraph{One word and joint generation.}
All two-adic source coordinates above are
finite products of continuous two-supported
powers of the actual $g_i,w_j$.
They therefore give one element of
$\widehat F'$ while clearing the full ternary
source image.  The three-supported
construction in E.37 does the converse.
The two corrections can be multiplied:
each is one in every class-two quotient at
the other's prime.  This realizes any
prescribed pair of full two-adic and
three-adic fibres over the same retained
source point.

The rank-ten tests imply that each finite
two-adic boundary map is an automorphism
of its actual $P_{2^n}$ quotient, by the
finite geometric Frattini criterion.
It descends to this quotient because it
preserves $H$ and the characteristic
subgroup $\gamma_3(H)H^{2^n}$.
The same is true on the three-adic factors.
Their common boundary quotient therefore
has an automorphism: its kernel is the
intersection of these preserved kernels,
and injectivity can be checked in every
factor.
For the joint source with $\mathscr D$,
apply the exact lower-central-stable-term
argument in E.36.  The $\mathscr D$ image
kills $\gamma_5(N_B)$, while a sufficiently
deep stable term of the common boundary
$N_B$ image occurs as a pure boundary
factor in the corrected image.
After quotienting it, the joint $N_B$
image is finite nilpotent with primes
two and three.  Both Frattini quotients
have rank $97$, forced by the retained
map through $(\mathbb Z/6)^{97}$.
The inherited $G_6$ automorphism then
gives surjectivity, and the quotient by
$N_B$ gives the whole common source.
This proves the generation assertion
without assuming kernel invariance for
the fixed Ferrari source.

The new exact certificates are
\texttt{bc\_fermat\_class\_two\_2adic\_source.py},
\texttt{bc\_fermat\_class\_two\_integral\_family.py},
\texttt{bc\_fermat\_class\_two\_2adic\_norms.py},
\texttt{bc\_fermat\_class\_two\_2adic\_fibres.py},
\texttt{bc\_fermat\_class\_two\_2adic\_words.py},
and \texttt{bc\_fermat\_class\_two\_2adic\_tower.py},
with their executed outputs and JSON/NPZ
data.  They retain the full source
presentations, all $432$ literal origins,
every finite P8 comparison and the
individual nonlifting word.  All finish
with \texttt{PASS}.
\end{proof}
